% bab3.tex — dari BAB_III_PENUTUP_POLISHED.md mulai ## A.

\section*{BAB III --- PENUTUP}
\addcontentsline{toc}{section}{BAB III --- PENUTUP}

\subsection*{A. Kesimpulan}

Pertama, taksonomi hadis merumuskan lima tahapan: pembentukan multi-saluran kenabian, penyebaran sahabat, pembakuan transmisi tabi'in, pembukuan kanonik ulama klasik, dan keragaman metodologis modern (Tanjung \& Fadhlurrahman, 2025; Musadad, 2026). Pascakeguncangan \textit{al-fitnah al-kubra}, keharusan \textit{sanad} dan semangat \textit{tathabbut} dipaksakan oleh keadaan; \textit{rihlah} kemudian mengukuhkannya sebagai kelaziman teknis hingga program \textit{tadwin} masa Umar bin Abdul Aziz (Tanjung \& Fadhlurrahman, 2025; Razak et al., 2025; Musadad, 2026). Berikutnya, penghimpunan \textit{Kutub al-Sittah} mematangkan kodifikasi, sedangkan trikotomi sahih-\textit{hasan}-\textit{dhaif} diberlakukan sebagai ukuran keotentikan berjenjang yang mengikat sarjana hingga kini. Rumusan masalah kesatu pun terjawab: tiap zaman memformalkan semangat penjagaan pendahulu ke perangkat yang kian baku --- dari ingatan-terapan ke kanon dan pangkalan data digital.

Kedua, \textit{hifz} ditegakkan sebagai basis kultural, \textit{kitabah} menopang ingatan, hingga \textit{tadwin} bersponsor penguasa menjadi muaranya (Rashwan, 2024; Lutfianto et al., 2024). Lewat jejaring guru-murid pengendali \textit{hifz}, \textit{sanad} berlandaskan \textit{'adalah} dan \textit{dabt} pun dilahirkan; sengketa pelarangan-perizinan ditengahi --- dari \textit{sadd al-dzari'ah} ke \textit{rukhsah} (Jaiyeoba \& Osmani, 2024; Watukila, 2024) --- sebagaimana dibuktikan \textit{Sahifah al-Sadiqah}. Penyebaran sahabat ke pelbagai wilayah menghasilkan khazanah kedaerahan yang tersaring lewat \textit{tathabbut} Khulafaur Rasyidin, lalu dirapikan melalui program \textit{tadwin} yang berpuncak pada \textit{al-Muwatta'} dan pola kanonik al-Bukhari sebagai kulminasinya (Gil, 2024; Obeid \& Hussein, 2023). Terjawablah rumusan kedua: penghimpunan berlangsung bertahap-berlapis --- khazanah tulis menyangga kelisanan; prakarsa perseorangan ditata otoritas publik hingga menjelma kanon bersama.

Ketiga, empat gugus pendorong pemalsuan teridentifikasi --- politis, teologis-sektarian, destruktif-infiltratif, devosional sosio-ekonomis --- dilacak melalui pemeriksaan rangkap \textit{sanad}-\textit{matan} berpasangan (Wasman et al., 2023; Muhamed Ali et al., 2015). Dorongan politis tumbuh dari kegentingan legitimasi pasca-\textit{fitnah} lewat \textit{afterlife} khilafah tiga dasawarsa; sementara unsur teologis, \textit{zindiq}, dongeng \textit{qussas}, dan anjuran \textit{targhib-tarhib} melanggengkan \textit{maudu'} (Liew, 2025; Su, 2021; Wasman et al., 2023; Noor, 2023). Pendeteksian memadukan uji \textit{sanad} --- keutuhan periwayat, ketersambungan rantai, terbuktinya \textit{liqa'} --- dengan telaah \textit{matan}: keselarasan Qur'an, dukungan riwayat kukuh, kesesuaian historis; vonis dijatuhkan bila indikator (Muhamed Ali et al., 2015; Calgan, 2024) terpenuhi secara kumulatif-kontekstual. Dengan demikian, separuh rumusan ketiga terjawab: pemalsuan berwatak krisis otoritas struktural, sehingga pendeteksiannya mensyaratkan kritik \textit{sanad-cum-matan} yang terlembaga dan berkesinambungan.

Keempat, penyelamatan disusun dalam tiga lapis: \textit{takhrij}, imunitas \textit{jarh-ta'dil}, dan penghimpunan \textit{maudu'at} sebagai gudang memori penangkal disinformasi (Muhamed Ali et al., 2015; Noor, 2023). \textit{Takhrij} kemudian ditransformasikan menjadi ekstraksi fitur komputasional plus pemetaan graf semantik; sedangkan \textit{jarh-ta'dil}, rintisan 'Ali b. al-Madini yang diwarisi Ibn al-Salah--Mamluk, disahkan lewat kajian jejaring atas \textit{Sahih Bukhari} berpendekatan jaringan (Wasman et al., 2023; Kamran et al., 2023; Su, 2021; Alam \& Schneider, 2020). Arsip \textit{maudu'at} rintisan Ibn al-Jauzi hingga al-Albani dialihfungsikan menjadi materi latih algoritme deteksi; telaah Ibn Kasir atas \textit{Sahihayn} menunjukkan tiada wilayah luput dari penilaian (Noor, 2023; Calgan, 2024). Separuh akhir rumusan ketiga pun menyimpulkan: kekebalan epistemologis hadis disangga oleh kesediaan diuji ulang lewat mekanisme kelembagaan yang ajek antar-generasi.

Kelima, makalah ini merangkai telaah historis-filologis klasik dengan pendekatan komputasional-digital --- meramu benang periodisasi, transmisi, dan pemalsuan secara bersintesis (Abdulrahman, 2024; Bin Jamil, 2024). Peranti \textit{pretrained}, \textit{AraBERT}, dan graf-\textit{transformer} diposisikan meneruskan \textit{tathabbut}, \textit{rihlah}, \textit{jarh-ta'dil} --- bukan menggantikan ulama --- sejalan dengan kerangka \textit{maqasid} sebagaimana dirumuskan (Bin Jamil, 2024; Abdulrahman, 2024). Urgensinya berpangkal pada konten viral yang hampa validitas dan \textit{ChatGPT} sebagai rujukan belajar, sehingga mendesak pengalihdigitalan \textit{takhrij} serta literasi hadis yang kritis-klasik (Abdullah et al., 2026; Syukur et al., 2026). Karenanya ilmu hadis tampil sebagai \textit{living hadith} yang lentur --- membentang dari \textit{sahifah} klasik ke pangkalan data masa kini (Abdulrahman, 2024; Matin Bin Salman, 2025).

\subsection*{B. Saran}

\begin{enumerate}[leftmargin=1.27cm,itemsep=2pt,parsep=2pt]
\item Perbandingan lintas mazhab-antargenerasi soal tafsir larangan penulisan layak diprioritaskan, sebab rekonsiliasi kronologis disandarkan (Lutfianto et al., 2024; Watukila, 2024) pada kepustakaan Sunni, 'Ajjaj al-Khatib, dan telaah Azami. Kepustakaan Syiah yang diperhadapkan dengan riset orientalis akan menguji daya laku umum kesepakatan \textit{nasakh} pelarangan lewat pembolehan penulisan menyeluruh.

\item Tolok ukur Arab-Indonesia untuk pemeriksaan digital memadukan petanda \textit{sanad} (kesinambungan dan sentralitas jaringan) dengan petanda \textit{matan} (gaya, sentimen, koherensi) dalam penilaian terbuka hasil kesepakatan (Mohamed \& Sarwar, 2022; Alghamdi et al., 2025). Berikutnya, korpus \textit{maudu'at} beranotasi \textit{'illah} dialihbentuk secara digital sebagai bahan latih penimbang setara model \textit{pretrained} dan \textit{AraBERT} secara berdampingan (Shaaban et al., 2026; Tarmom et al., 2022).

\item Literasi keagamaan kaum muda seyogianya membekalkan \textit{takhrij} digital secara bertahap --- dari penelusuran daring sampai pengecekan \textit{sanad}-\textit{matan} --- untuk menanggapi \textit{ChatGPT} sebagai bahan belajar dan viralnya hadis tanpa keabsahan (Syukur et al., 2026; Abdullah et al., 2026). Riset aksi (\textit{action research}) literasi hadis digital yang telah dirintis selayaknya diperbanyak di jenjang menengah dan dakwah dalam jaringan (Supriyadi et al., 2020; Abdulrahman, 2024).

\item Kritik \textit{sanad}-\textit{matan} dalam kurikulum keislaman --- dari pesantren hingga kampus --- seyogianya mengangkat \textit{jarh-ta'dil}, \textit{takhrij}, dan telaah kontekstual menjadi kecakapan yang melampaui hafalan semata (Muhamed Ali et al., 2015; Wasman et al., 2023). Kebangkitan Mamluk lewat mediasi \textit{Muqaddimah} Ibn al-Salah menunjukkan bahwa pembakuan kurikulum selalu lahir untuk menjawab kegentingan; teladan ini pantas ditempuh ulang saat ini (Noor, 2023; Bin Jamil, 2024).

\item Kajian prosopografis jejaring periwayat Hijaz, Irak, Syam, dan Mesir mengawinkan telaah jejaring dengan graf pengetahuan berbasis \textit{transformer} untuk membidik anomali transmisi berskala luas (Alam \& Schneider, 2020; Mosa, 2025). Lewat pendekatan itu, ketaksaan identitas narator dapat disempurnakan, dan dugaan kehati-hatian berlebih pada wilayah tertentu dapat diuji memakai data terbuka-tereplikasi menurut telaah (Kamran et al., 2023; Su, 2021).

\end{enumerate}
